This chapter covers all 38 MCCME records numbered 92--129 and all 68 supplied scans, P000140--P000207. The catalogue order is preserved even when dates are not monotone. Its overlapping page ranges identify records; each occurrence instead points to the pages containing that actual statement. The year heading, prospective plan, and Parshin heading remain present without invented mathematical content. The outcome is a critical occurrence index of 277 items, 21 indexed problems (one continuing an earlier problem identifier), 108 correction or qualification entries, and 78 checked bibliographic entries. Status assessments are dated 9 November 2026.

\begin{center}
\begin{tabular}{lrrrrrr}
\toprule
Class & A & B & C & D & E & F\\
\midrule
Items & 22 & 45 & 3 & 0 & 185 & 22\\
\bottomrule
\end{tabular}
\end{center}
Here A denotes an authorial publication soon after the talk, B subsequent development by others, C a precisely delimited partly developed question, D an undeveloped identifiable item, E prior mathematics or exposition, and F failure of exact identification. These are historical classifications, not correctness scores. An item may be A and false as written; a reconstruction may be correct without proving what the speaker intended. The separate assessment fields preserve these distinctions.

Three lines dominate. Bernstein’s equivariant K-theory and elliptic-family constructions connect representation rings, completion, and index bundles. The recorded claim that every representation ring embeds in its completion fails for a disconnected compact group; a concrete cyclic-group counterexample is supplied. Minlos’s lectures move from canonical models and factorial correlations to thermodynamic quantities and limiting distributions. The later published lectures explicitly trace their 1965 and 1966 origins \citep[Introduction]{minlos1968lectures}. The year also follows the expansion of singularity theory: jets, Boardman symbols, intrinsic derivatives, preparation, and real local algebras. Arnold’s two November lectures are singularity-theory lectures despite the catalogue’s hydrodynamics label \citep{arnold1968singularities,boardman1967singularities}.

The strongest publication matches include Kazhdan’s property-\(T\) announcement, Shilov’s regular Hilbert-space analysis, Seeley’s complex-power kernel statements, and Vishik--Eskin’s family-index criterion for completion of elliptic boundary systems \citep{kazhdan1967dual,shilov1967hilbert1,seeley1968powers,vishikeskin1967normally}. In the last case the manuscript received in April 1967 resolves the March question within its stated setting; it is not counted as an orphan problem. Dates of talks, receipt, preprint circulation and printing are kept separate whenever the evidence distinguishes them.

The three C items concern the algebraic topological classification of Anosov diffeomorphisms, the arbitrary-smooth-metric remainder of Helgason’s exact geodesic-sphere support question, and the smooth type of equal-prime-period geodesic manifolds. Each is matched with its own hypotheses; a geodesic-ray theorem, an analyticity theorem, or an entire-Grauert-tube theorem is not substituted for the recorded problem \citep[Conjecture 2.3]{spatzier2016rigidity}\citep[Problem 4.1]{quinto2008helgason}\citep[Main Theorem]{lisong2026zoll}. No D item is justified, and none of these C items yielded a plausible bounded research advance in this run.

The full mathematical annotation therefore audits Minlos’s variance argument. It proves an explicit conditional variance estimate, an \(L^1\) truncation conclusion, and counterexamples to finite variance under a bare \(L^1\) hypothesis. It also identifies the exact canonical volume correction and the legitimate dependence of the effective volume on correlation order. These are complete, independently checked editorial derivations, with code and outputs, and carry no claim of a new Gibbs-state theorem. The larger thermodynamic programme’s undefined exceptional parameters are retained as a formulation limitation (F).

Continuity with earlier chapters includes lattice rigidity and cohomology (1964-8 and 1964-24), family indices (1964-15.4), the Maslov thread (1964-21.3), and the pseudodifferential calculus of 1965. The perturbed complex square reuses GP-1965-03 and its earlier occurrence 1965-72.5. The corrections distinguish changed readings, translations, unresolved ink, provenance, and faithful mathematical errors: the Boardman sign and positive-sequence convention, the elliptic local-algebra monomial, entropy normalization, and the heat-kernel factor four are consequential examples. Every source was opened, but a verified link does not mean that every proof was read; the item assessments and bibliography annotations state the actual verification level.
